\section{Introduction}
\label{sec:intro}

The Brazilian equity market is among the most volatile in the emerging-market
universe, where returns are shaped not only by firm-level fundamentals but also by
sharp swings in the domestic monetary policy cycle.
Between 2021 and 2023, the Banco Central do Brasil raised the benchmark SELIC rate by
523 basis points---from 2.00\% to 13.75\%---the most aggressive tightening cycle in
two decades, triggered by a combination of post-pandemic supply shocks and persistent
inflationary pressure.
This episode illustrates a structural feature of the Brazilian market: asset-return
dynamics shift substantially across interest-rate regimes, with aggregate Sharpe ratios
varying from $-0.23$ during rising-rate periods to $+0.88$ during stable-rate periods
in our sample.
Classical portfolio optimization frameworks, however, are largely regime-agnostic.
Mean-variance optimization \citep{Markowitz1952} and its variants estimate a single
covariance structure over the entire sample, ignoring the non-stationarity induced by
monetary-policy transitions.
Rule-based approaches such as the $1/N$ strategy \citep{DeMiguel2007} and Risk Parity
\citep{Qian2011} are robust to estimation error but equally insensitive to prevailing
macroeconomic conditions.

A natural extension is to use reinforcement learning (RL), which can, in principle,
learn a state-dependent policy that adapts asset weights to the current economic
environment.
Online RL agents trained on rolling windows of live market data have been applied to
portfolio management with promising results \citep{Jiang2017}, but this paradigm
introduces a subtle methodological problem: the agent can implicitly learn from
future information during training because the reward is computed on the same data
stream that updates the policy, creating a look-ahead bias that inflates reported
performance.
Offline RL \citep{Levine2020} resolves this by training entirely on a fixed, historical
dataset---the offline dataset---and deploying the learned policy only on held-out
data.
Conservative Q-Learning (CQL) \citep{Kumar2020} is a state-of-the-art offline RL
algorithm that avoids the overestimation of Q-values for out-of-distribution actions
by adding a conservatism penalty to the Bellman loss, making it particularly well
suited to the noisy, non-stationary nature of financial time series.

Despite the growing body of work applying RL to asset allocation, three important
gaps remain in the literature.
First, most studies focus on developed markets; evidence from the Brazilian market,
with its distinctive macro-financial dynamics, is scarce.
Second, the predominant reward signal is the Sharpe ratio, which treats upside and
downside volatility symmetrically and can incentivize the agent to accept large tail
losses.
Third, offline RL methods have rarely been combined with explicit macroeconomic
regime information in the state representation, leaving the potential of
regime-conditional portfolio decisions largely unexplored.

This paper addresses all three gaps.
Our specific contributions are as follows:

\begin{enumerate}
  \item \textbf{Regime-aware state representation.}
        We construct a 225-dimensional state vector that combines asset-level
        return statistics with macroeconomic variables---SELIC regime indicators
        (one-hot encoded as rising, stable, or falling), CDS spreads, credit
        spreads, and Brazilian implied volatility---enabling the agent to condition
        portfolio decisions on the prevailing monetary-policy environment.

  \item \textbf{CVaR reward function.}
        We replace the standard Sharpe-based reward with a composite signal that
        penalizes tail risk via the Conditional Value-at-Risk at the 95\%
        confidence level while retaining a linear expected-return component
        weighted by a risk-aversion parameter $\lambda = 0.5$.
        This incentivizes the agent to control downside risk explicitly, a
        property especially valuable in the fat-tailed return distributions
        observed in emerging markets.

  \item \textbf{Rigorous out-of-sample evaluation.}
        We employ a strict walk-forward protocol with 35 quarterly folds, training
        on 252-day windows and evaluating on 63-day windows, ensuring zero
        look-ahead bias.
        We benchmark the CQL agent against Equal-Weight, Minimum-Variance, and
        Risk-Parity portfolios and conduct an ablation study that isolates the
        contributions of the offline paradigm, the CVaR reward, and the macro
        state variables.
\end{enumerate}

The remainder of the paper is organized as follows.
Section~\ref{sec:litreview} reviews the related literature on classical portfolio
optimization, reinforcement learning in finance, offline RL, CVaR-based objectives,
and macroeconomic regime detection.
Section~\ref{sec:data} describes the dataset, data sources, and regime
classification.
Section~\ref{sec:methodology} presents the MDP formulation, the CQL algorithm,
the reward function, and the evaluation protocol.
Section~\ref{sec:results} reports out-of-sample results and the ablation study.
Section~\ref{sec:discussion} interprets the findings and discusses limitations.
Section~\ref{sec:conclusion} concludes.
